% !TEX program = xelatex
\documentclass{英语试卷}
\begin{document}
\examtitle{2024-2025 学年福州市高三年级第四次质量检测}{英语}
\par 英语学科 试题卷\par
\begin{notice}
\par 1. 答卷前，考生务必将自己的姓名、准考证号填写在答题卡上，\par
\par 2. 回答选择题时，选出每小题答案后，用铅笔把答题卡上对应题目的答案标号涂黑。如需改动，用橡皮擦干净后，再选涂其他答案标号。回答非选择题时，将答案写在答题卡上，写在本试卷上无效。\par
\par 3. 考试结束后，将本试卷和答题卡一并交回。\par
\end{notice}
\parttitle{第一部分 听力（共两节，满分 30 分）}
\par 做题时，先将答案标在试卷上，录音内容结束后，你将有两分钟的时间将试卷上的答案转涂到答题卡上。\par
\sectitle{第一节（共5小题：每小题1.5分，满分7.5分）}
\par 听下面 5 段对话。每段对话后有一个小题，从题中所给的 A、B、C 三个选项中选出最佳选项。听完每段对话后，你都有 10 秒钟的时间来回答有关小题和阅读下一小题。每段对话仅读一遍。\par
\question{1. What is Mr. Smith like?}
\begin{choices}(3)
  \item Strict.
  \item Professional.
  \item Controlling.
\end{choices}
\question{2. What will the man probably do next?}
\begin{choices}(3)
  \item Call Henry for help.
  \item Buy a new machine
  \item Repair the machine himself.
\end{choices}
\question{3. What are the speakers mainly talking about?}
\begin{choices}(3)
  \item A nature film
  \item Antarctic weather.
  \item Penguin behavior.
\end{choices}
\question{4. What will the woman do this evening?}
\begin{choices}(3)
  \item Work overtime.
  \item Tidy up her place.
  \item Hang out with the man.
\end{choices}
\question{5. Why does the man sell his car?}
\begin{choices}(3)
  \item It's no longer reliable.
  \item He does not drive anymore.
  \item He wants to lower travel costs.
\end{choices}
\sectitle{第二节（共15小题：每小题1.5分，满分22.5分）}
\par 听下面 5 段对话或独白。每段对话或独白后有几个小题，从题中所给的 A、B、C 三个选项中选出最佳选项。听每段对话或独白前，你将有时间阅读各个小题，每小题 5 秒钟；听完后，各小题将给出 5 秒钟的作答时间。每段对话或独白读两遍。\par
\question{6. What is Lisa doing recently?}
\begin{choices}(3)
  \item Writing school papers.
  \item Making a diet schedule.
  \item Studying for final exams.
\end{choices}
\question{7. What does the man suggest Lisa do?}
\begin{choices}(3)
  \item Redo work tomorrow.
  \item Switch to healthier meals.
  \item Create a balanced study plan.
\end{choices}
\par 听第 7 段材料，回答第 8 至 10 题。\par
\question{8. What service does PictureRem provide?}
\begin{choices}(3)
  \item Repairing damaged photos.
  \item Decorating old photos with covers.
  \item Turning photos into digital copies.
\end{choices}
\question{9. What does the man ask the woman to do?}
\begin{choices}(3)
  \item Bring photos with covers.
  \item Bring photos out of albums.
  \item Bring photos fewer than 100.
\end{choices}
\question{10. How much should the woman pay in total?}
\begin{choices}(3)
  \item \$ 50.
  \item \$ 85.
  \item \$ 90.
\end{choices}
\par 听第 8 段材料，回答第 11 至 13 题，\par
\question{11. What is the film Feeling from Mountain and Water about?}
\begin{choices}(3)
  \item A boy taught to play the guzheng.
  \item A boy exploring the magical nature.
  \item A musician traveling through mountains
\end{choices}
\question{12. What led to the decline of ink-wash art after the 1980s?}
\begin{choices}(3)
  \item High cost of producing.
  \item Loss of visual symphony.
  \item Lack of audience interest.
\end{choices}
\question{13. How does the ink-wash art come back to life?}
\begin{choices}(3)
  \item By following traditional methods.
  \item By integrating with other styles.
  \item By using digital techniques.
\end{choices}
\par 听第 9 段材料，回答第 14 至 17 题，\par
\question{14. What is Mr. Chen doing?}
\begin{choices}(3)
  \item Hosting a charity event.
  \item Giving a public speech.
  \item Participating in an interview.
\end{choices}
\question{15. What do we learn about Mr. Chen's childhood?}
\begin{choices}(3)
  \item He often skipped classes.
  \item He lived in a wealthy family
  \item His parents valued education.
\end{choices}
\question{16. What did Mr. Chen major in at college?}
\begin{choices}(3)
  \item Primary education.
  \item Ocean engineering.
  \item Environmental science
\end{choices}
\question{17. How does Mr. Chen feel about Lily?}
\begin{choices}(3)
  \item Grateful.
  \item Surprised.
  \item Disappointed.
\end{choices}
\par 听第 10 段材料，回答第 18 至 20 题。\par
\question{18. Where did jazz first appear?}
\begin{choices}(3)
  \item In Africa.
  \item In America.
  \item In New Zealand.
\end{choices}
\question{19. How did Louis Armstrong influence jazz?}
\begin{choices}(3)
  \item He changed its musical style.
  \item He improved live stage performances.
  \item He invented new musical instruments.
\end{choices}
\question{20. What impact did jazz have beyond music?}
\begin{choices}(3)
  \item It standardized thinking ways.
  \item It increased dance types.
  \item It encouraged freedom.
\end{choices}
\parttitle{第二部分 阅读理解（共两节，满分 50 分）}
\sectitle{第一节（共15小题：每小题2.5分，满分37.5分）}
\par 阅读下列短文，从每题所给的 A、B、C 和 D 四个选项中，选出最佳答案，\par
\passagelabel{A}
\par Genspace, a community biolab in Brooklyn, NY, offers the Biorocket Internship (实习) to high school students passionate about biotechnology. This program is designed to provide hands-on experience in a real laboratory setting, allowing students to work on cutting-edge projects and gain invaluable skills that can set them apart in their college applications.\par
\par What does a typical day as a Biorocket intern look like?\par
\par A typical day in the Biorocket Internship starts with a brief lecture or workshop on a specific topic, followed by real-world lab work. Students might spend the morning learning about DNA extraction (提取) techniques and then apply that knowledge in the lab during the afternoon. There are also regular meetings with guest speakers who are experts in various fields of biotechnology, providing students with insights into potential career paths and the latest industry trends.\par
\par Who can apply?\par
\par Students who are at least 16 years old and currently attending a NYC public school within 45 minutes of Genspace are qualified to apply. By applying, you're showing commitment to participating in the program from February to August.\par
\par Application Requirements\par
\par A personal profile\par
\par Academic information\par
\par Responses to short essay questions\par
\par Contact information for their parents and one teacher reference\par
\par Note: Materials should be submitted by the deadline, usually in early January. The selection committee reviews these materials to assess the applicant's academic achievements, commitment to the field, and potential for growth. Early preparation and thorough attention to the application components can significantly enhance the chances of being selected for this distinguished internship.\par
\question{21. What is the aim of Genspace?}
\begin{choices}(2)
  \item To fund biological students.
  \item To provide lab experience.
  \item To observe science learning.
  \item To present a rocket project.
\end{choices}
\question{22. What are the Biorocket interns expected to do?}
\begin{choices}(2)
  \item Seek different career paths.
  \item Put bio-knowledge into practice.
  \item Host meetings with experts.
  \item Give lectures on related topics.
\end{choices}
\question{23. What does the selection committee assess?}
\begin{choices}(3)
  \item Physical condition. 3. Family history.
  \item Learning potential.
  \item Field research.
\end{choices}
\passagelabel{B}
\par Kristy Gordon's artistic voice didn't emerge overnight—it was a year-long process of unlearning perfectionism and rediscovering her creativity. Like many artists, she began by burying her intuition (直觉) under academic expectations. Through strict training in classical techniques, she mastered the rules of "good" painting—only to realize those very rules had become a box.\par
\par The turning point came after completing her MFA at the New York Academy of Art. Free from the pressure of peer review, Gordon gave herself permission to experiment without judgment. Gordon made a brave decision to step away from public exhibitions for an entire year, which gave her the freedom to explore without worrying about critical reception or commercial success. Without the approaching deadline of an upcoming exhibition, she found herself painting more freely, trying techniques she'd previously considered "too risky" for public display.\par
\par Gordon rediscovered the pure joy of creating for herself alone. The year away from the spotlight allowed her to develop a body of work that was truly hers, rather than shaped by market trends or academic feedback. As she later recalled, "That year of painting in private was like giving myself permission to be a beginner again. Without an audience, I could make all the 'mistakes' that eventually led to my artistic breakthroughs."\par
\par The strategy paid off—when she eventually returned to exhibiting, her work showed a new level of confidence and originality that had been formed in that year.\par
\par Gordon's journey reveals an insightful truth: An artist's voice isn't something to be found, but something to be uncovered. It requires both courage to silence external expectations and patience to hear your own inner creative whisper. As she reflects, "Your most authentic (真实的) work emerges when you stop trying to make art, and start allowing art to make you."\par
\question{24. Which best describes Gordon's previous paintings?}
\begin{choices}(2)
  \item Impressionistic.
  \item Conventional.
  \item Abstract.
  \item Expressive.
\end{choices}
\question{25. How does Gordon start to find her voice?}
\begin{choices}(2)
  \item She posted paintings online.
  \item She gave up risky techniques.
  \item She took diverse art training.
  \item She kept a free creative space.
\end{choices}
\question{26. Which statement best explains Gordon's comment about "mistakes"?}
\begin{choices}(2)
  \item Unrestricted creation brings innovation.
  \item Original art is often rejected by galleries.
  \item Market trends hold back art development.
  \item Beginners should hide unsatisfying works.
\end{choices}
\question{27. What is Gordon's understanding of art?}
\begin{choices}(2)
  \item It is a competition for recognition.
  \item It should follow strict regulations.
  \item It is an exploration of inner self.
  \item It should challenge expectations.
\end{choices}
\passagelabel{C}
\par Large language models (LLMs) seem better at making people feel seen and heard than humans. This phenomenon, called "LLMpathy", is both surprising and controversial. Some argue that computers can't truly empathize (共情) for lack of emotion. Others are alarmed by how readily people are trading human connection for digital ones. But beyond these concerns, chatbots might offer something more practical. If they are beating us at empathy, shouldn't we try to learn what they are doing, right?\par
\par Researchers initially wondered if AI advantage lay in its unlimited attention. But that doesn't seem to explain it. A 2024 research in Harvard Business School revealed that people's expressing of empathy, despite increased commitments, still fell short compared to those by ChatGPT.\par
\par Chatbots' success may come from avoiding all-too-human mistakes. Chatbots' responses stood out mainly because they focused on acknowledging feelings, unlike humans who often shared a seemingly related experience to offer solutions, unintentionally damaging the hope of being truly heard.\par
\par Chatbots avoid these pitfalls. They focus entirely on the speaker. More than humans, chatbots paraphrase, acknowledge and justify how people might feel. When people adopt similar strategies, their connections strengthen.\par
\par It bears noting that the AI advantage in empathetic conversations has limits. Its "paraphrase, acknowledge, follow up" may feel warm the first time, but dull the second. Most research tests people's interactions with Chatbots just once, so AI's edge might fade in longer, repetitive chats.\par
\par Chatbots might be helpful, but they still can't feel or truly care. The demand for AI therapists (理疗师) may be growing, but many people still prefer human support. Anyone who has repeated "agent" at a customer-service robot knows the feeling of desperately wanting a real person on the line. Some of the shortcomings of human connection are also, in fact, features. But the fact that we often must earn human empathy, and that it comes from limited beings who devote themselves to being there for us, is part of its beauty.\par
\question{28. Why does AI often outperform humans at showing empathy?}
\begin{choices}(3)
  \item It has unlimited patience. 3. It shares related experience
  \item It focuses on people's feeings
  \item It prioritizes solving problems.
\end{choices}
\question{29. What does the underlined word "pitfalls" in paragraph 4 probably mean?}
\begin{choices}(2)
  \item Responses.
  \item Shortcomings.
  \item Experiences.
  \item Feelings
\end{choices}
\question{30. What can we learn from the last paragraph?}
\begin{choices}(2)
  \item Human care is irreplaceable.
  \item In-person chats have limits.
  \item Service lines benefit human.
  \item Emotion is enhanced by AI.
\end{choices}
\question{31. Which of the following is a suitable title for the text?}
\begin{choices}(2)
  \item Chatbots: The New Emotional Support
  \item The Double-Edged Sword of Chatbots
  \item The Rising Demand for AI Therapists
  \item Learn to Be a Listener like Chatbots
\end{choices}
\passagelabel{D}
\par Soratobu Medaka, a children's book, tells the story of tiny medaka fish riding on a bird to a distant pond. It has delighted children, but now it's revealed to be true.\par
\par Researchers have long struggled to explain how fish turn up in lakes far from other bodies of water, a mystery even Charles Darwin noted in 1859, when he suggested that aquatic (水生的) hatchlings might stick to the feet of birds, though his bright idea remained a theory.\par
\par In 2019, Yao Akifumi, a biology student then, wondered if Darwin might have been right and further wondered if fish-hunting birds might transport egg-bearing aquatic plants stuck to their feet. To test this, he set up two ponds one metre apart, one containing 36 plants and the other without, both equipped with camera traps. After six months, some plants had switched ponds, and one camera captured a bird flying away with a plant on its foot. Next, Yao aimed to see if fish eggs could weather the journey being carried out of water. He attached tiny medaka eggs to aquatic plants and removed them for zero to 24 hours—at favorable hatching conditions—before returning them.\par
\par In a paper in Science of Nature, Mr. Yao reports that the medaka eggs can successfully hatch after up to 18 hours out of the water. Eggs without a plant leaf to attach to rarely survived a few hours. The big drop in survival was at 16.3 hours, when half the eggs became lifeless.\par
\par Whether any egg could survive that long, given that a bird's average flight-speed is around 39kph, is doubtful, says Mr. Yao, since wind exposure during such a long journey would dry them out.\par
\par Shorter journeys, however, are certainly possible. And since many other medaka-like species are found throughout Asia, the chances are high that hitchhiking (搭便车) explains how they have ended up in lakes with seemingly no connection to other water bodies. Yet again, it would seem, Darwin's theories turn out to be correct.\par
\question{32. Why does the author mention Darwin's idea?}
\begin{choices}(2)
  \item To introduce an academic term.
  \item To provide research background.
  \item To question a well-known theory.
  \item To explore new research methods.
\end{choices}
\question{33. What evidence did Yao find to support his assumption?}
\begin{choices}(2)
  \item Birds flying between ponds.
  \item Fish eggs surviving on plants.
  \item A plant growing in both ponds.
  \item A bird carrying a plant on its foot.
\end{choices}
\question{34. What may increase medaka eggs' survival?}
\begin{choices}(2)
  \item Birds' flying route.
  \item Wind exposure.
  \item Supporting leaves.
  \item Longer journeys.
\end{choices}
\question{35. What is the main idea of the text?}
\begin{choices}(2)
  \item Children's book inspires research.
  \item Origin of fish has been uncovered.
  \item Fish's hitch riding has been proven
  \item Eggs and plants have a close relation.
\end{choices}
\sectitle{第二节（共5小题；每小题2.5分，满分12.5分）}
\par 阅读下面短文，从短文后的选项中选出能填入空白处的最佳选项。选项中有两项为多余选项，\par
\par The great irony is that as we become increasingly "connected"—on social media, video calling, and messaging—we still feel increasingly lonely. And even though we may use technology to feel more connected, it may be exactly what's leading us to feel lonely.\par
\par What about you? \blanknum{36} If not, try some of these strategies to stop feeling lonely.\par
\par Talk to strangers. A growing body of research suggests that even seemingly worthless interactions with strangers—like chatting with a cashier—may be able to keep loneliness at bay by helping us feel more socially connected. So reach out to other human beings to say hello, or chat about whatever's on your mind. \blanknum{37}\par
\par Stop focusing so much on you. It's almost unavoidable in our modern technology-crazed world that we start to believe we don't have enough. Instead of focusing on what you can get, shift your focus to what you can give. You could sell T-shirts online to raise money for a good cause. You could ask friends to donate to a charity for your birthday. \blanknum{38} These small acts can make a big difference and help you reduce feelings of loneliness.\par
\par Connect in real life. Connecting in real life may not be as easy as it once was. We often default (默认) to using our smartphones—it's easier, and now it's culturally accepted. \blanknum{39} We do this by looking people in the eyes, listening, being mindful, and choosing not be distracted by our phones or other technologies.\par
\par \blanknum{40} It's important to practice self-sympathy when you fail at things. There is no need to put yourself down. That kind of attitude won't help you decrease loneliness. Instead, try talking to yourself in a way that is supportive and caring—and you'll be more likely to acknowledge mistakes in trying to decrease loneliness, and hopefully do better next time.\par
\begin{choices}(1)
  \item Be nice to yourself.
  \item Stop your negative thought cycles.
  \item Are you feeling socially connected?
  \item Do you find online interactions less satisfying?
  \item It helps you to feel more connected and less lonely.
  \item But we can decrease our loneliness if we build stronger personal interaction.
  \item By giving to others, you take the focus off yourself and do good at the same time.
\end{choices}
\parttitle{第三部分 语言知识运用（共两节，满分 30 分）}
\sectitle{第一节（共15小题；每小题1分，满分15分）}
\par 阅读下面短文，从短文后各题所给的 A、B、C 和 D 四个选项中，选出可以填入空白处的最佳选项，\par
\par Mariatou, a woman from a small village in Senegal, had always dreamed of building a better life for her family. However, her husband's illness and the lack of \blanknum{41} in her village made life extremely challenging. Cooking in the dark was a daily \blanknum{42} Children also had difficulty studying, \blanknum{43} kerosene (煤油) lamps that were inefficient.\par
\par When Barefoot College International \blanknum{44} a solar training program in her village, Mariatou was initially \blanknum{45} due to her inability to read and fear of the unknown. But driven by her love for her family and the \blanknum{46} need for electricity in her community, she decided to \blanknum{47} and sign up. She recalled, "I thought, 'You cannot always find all the \blanknum{48} in your village, sometimes you have to go out to explore more'." The training \blanknum{49} her life. She went from being a woman with limited opportunities to a skilled solar engineer, \blanknum{50} of setting up and maintaining solar products. This \blanknum{51} knowledge gave her the confidence to handle new challenges and spread her techniques to others.\par
\par Now, Mariatou is ready to bring light and \blanknum{52} to her village. She is grateful to Barefoot College International and all those who \blanknum{53} her journey. Her story proves the power of \blanknum{54} and access to resources. It shows that with determination and \blanknum{55}, even the wildest dreams can become reality. Mariatou's success not only benefits her family but also serves as an inspiration to others in her community.\par
\cz{41}{electricity}{water}{sunlight}{food}
\cz{42}{joy}{lesson}{target}{struggle}
\cz{43}{hanging up}{relying on}{turning down}{cleaning up}
\cz{44}{completed}{joined}{offered}{expanded}
\cz{45}{hesitant}{curious}{excited}{surprised}
\cz{46}{emotional}{urgent}{equal}{social}
\cz{47}{make a point}{take a chance}{have an excuse}{do a favour}
\cz{48}{solutions}{choices}{problems}{promises}
\cz{49}{saved}{enriched}{risked}{transformed}
\cz{50}{afraid}{aware}{capable}{sick}
\cz{51}{limited}{self-taught}{detailed}{new-found}
\cz{52}{hope}{business}{fame}{success}
\cz{53}{followed}{supported}{shared}{recorded}
\cz{54}{culture}{nature}{education}{friendship}
\cz{55}{patience}{gratitude}{cooperation}{opportunity}
\sectitle{第二节（共10小题；每小题1.5分，满分15分）}
\par 阅读下面短文，在空白处填入 1 个适当的单词或括号内单词的正确形式。\par
\par As an 80s-born individual coming from a rural north, I moved to the city and had a heavy diet, often finishing leftovers at home. \blanknum{56} (combine) with stressful urban life and work, I gradually piled on the pounds. At my \blanknum{57} (heavy), I tipped the scales at 98 kilograms, while my height measured 173cm. I decided that the time had come to effect \blanknum{58} change.\par
\par I was firmly convinced that the key lies in one's own self-awareness. To attain my weight-loss objective, I \blanknum{59} (primary) concentrated on regulating my diet, cut down on my carbohydrate intake and \blanknum{60} (have) more vegetables and lean proteins. Unsweetened black coffee also emerged \blanknum{61} a steady companion, which gave a boost to my metabolism (新陈代谢), contributing to my overall \blanknum{62} (well) and vitality significantly. Additionally, I made \blanknum{63} a fixed daily routine to take a walk, regardless of how busy my schedule was.\par
\par Through continuous effort, I eventually succeeded. Losing weight made me feel lighter, more energetic and healthier. I am filled with more energy \blanknum{64} (deal) with my daily work and life. So far, my health has witnessed a substantial improvement.\par
\par This weight-loss journey, \blanknum{65} I overcame numerous difficulties, has taught me a profound lesson: as long as we stay determined and take practical actions, we can achieve the goals that we pursue.\par
\parttitle{第四部分 写作（共两节，满分 40 分）}
\sectitle{第一节（满分15分）}
\par 假定你是李华，你的笔友 Tom 对你校课间 “特色体育活动” 很感兴趣，请你给他写一封邮件，内容包括：\par
\par (1) 体育活动内容;\par
\par (2) 同学们的感受。 注意：\par
\par (1) 写作词数应为 80 左右;\par
\par (2) 请按如下格式在答题纸的相应位置作答，\par
\par Dear Tom,\par
\par I'm glad to share with you the featured sports activities during break in our school.\par
\sectitle{第二节（满分25分）}
\par 阅读下面材料，根据其内容和所给段落开头语续写两段，使之构成一篇完整的短文。\par
\par As class monitor, I was used to organizing group work, not helping students alone. One cold October afternoon, Mrs. Jenkins asked me to stay after class. Her usual smile was gone. "Ethan Miller joined us last month," she said quietly. "His family is going through hard times. His parents work two jobs day and night to make ends meet. He's failing classes and... stays silent, avoiding everyone." She looked at me. "Will you help him? You might be the only one he'll talk to."\par
\par I froze, holding my books tightly. The final exams were six weeks away, and my own grades had stopped improving. Even after staying up late studying math, my test scores stayed the same. My parents' worried faces came to mind. Helping Ethan meant losing precious exam preparation time, but I couldn't say no. "I'll try," I said weakly, though I felt uneasy—what if this hurt us both?\par
\par Ethan was like a shadow in class. He sat at the back, head down, writing notes but never speaking. At lunch, he ate alone by the lockers. When I first greeted him, he only nodded. After three failed tries to chat, I grew impatient. "We'll prepare for final exams together after school," I said firmly, using my role as monitor. He stayed silent but didn't refuse.\par
\par Our first study meeting was awkward. I opened my textbooks at the library table while Ethan sat rigidly, staring at his old shoes. I repeated history facts and science notes like a machine; he copied them quietly. During our third session, I explained a chemistry problem but made a mistake. Ethan pointed at my paper. "This number is wrong," he said softly. I checked—he was right. "How did you know?" I asked, surprised. He paused. "I reviewed the chapter last night... didn't want to waste your time." For the first time, I saw the tiredness in his eyes—just like mine from late-night studying.\par
\par 注意：\par
\par (1) 续写词数应为 150 左右;\par
\par (2) 请按如下格式在答题卡的相应位置作答。\par
\writingpara{Paragraph 1: Slowly, things changed between us.}
\writinglines{8}
\writingpara{Paragraph 2: On results day, we were amazed by our scores.}
\writinglines{8}
\end{document}
